% !PRINTED 543
\weberpage{543}
%<<CONT
$\pi_i$, but not divisible by others, and this requirement can always
be satisfied according to \S.\,143, 3. Then by 2. we can determine
the integer $\tau$ so that
\[
\alpha\xi + \beta_1\tau = 1
\]
holds, and if we set $\tau\eta_1 = \eta$, $\tau\zeta_1 = \zeta, \ldots$,
we obtain the theorem:

3. \gesperrt{If $\alpha, \beta, \gamma, \ldots$ are numbers in
$\mathfrak{o}$ without a common divisor, then other numbers
$\xi, \eta, \zeta, \ldots$ in $\mathfrak{o}$ can be determined so that}
\weq{4}{\alpha\xi + \beta\eta + \gamma\zeta + \cdots = 1}
\gesperrt{holds.}

From this we can further infer the following theorem:

4. \gesperrt{If $\alpha, \beta, \gamma, \ldots$ are arbitrary numbers
in $\mathfrak{o}$, and $\mu$ is a number in $\mathfrak{o}$ divisible
by the greatest common divisor of all these numbers, then the numbers
$\xi, \eta, \zeta, \ldots$ in $\mathfrak{o}$ can be determined so that}
\weq{5}{\mu = \alpha\xi + \beta\eta + \gamma\zeta + \cdots}
\gesperrt{holds.}

For if, according to \S.\,139, we take the greatest common divisor of
the numbers $\alpha, \beta, \gamma, \ldots$ in the form
\[
\delta = \alpha u + \beta v + \gamma w \ldots
\]
in which $u, v, w, \ldots$ are variables, then by assumption there is
an integral functional $\omega$ such that $\mu = \delta \omega$. We
represent the functional $\omega$ as a quotient of two integral
functions $\varphi : \varepsilon$, of which $\varepsilon$ is a unit,
and consequently $\varphi$ a function with integral coefficients, and
so obtain
\weq{6}{\varepsilon\mu = (\alpha u + \beta v + \gamma w + \cdots)\varphi.}
If one arranges both sides of this equation according to the variables
occurring in them, one obtains, when $\varepsilon_1, \varepsilon_2, \ldots$
are the coefficients in $\varepsilon$, a system of equations of the
following form:
\weq{7}{\begin{aligned}
\varepsilon_1\mu &= \alpha\xi_1 + \beta\eta_1 + \gamma\zeta_1 + \cdots \\
\varepsilon_2\mu &= \alpha\xi_2 + \beta\eta_2 + \gamma\zeta_2 + \cdots \\
& \text{.\,.\,.\,.\,.\,.\,.\,.\,.\,.\,.\,.\,.\,.\,.\,.\,.\,.\,.\,.\,.\,.\,.,}
\end{aligned}}
in which the $\xi_i, \eta_i, \zeta_i, \ldots$ are numbers in
$\mathfrak{o}$. But now the numbers $\varepsilon_1, \varepsilon_2, \ldots$,
as coefficients of a unit, have no common
%CONT>>
